\subsection{Continuum Balance Laws as Motivation for 1D Models}
\label{subsec:governing-equations-model-hierarchy}

The full-order starting point is a three-dimensional incompressible flow model on
a lumen with inlet, outlet, and wall boundary portions. A compact reading is
sufficient for this report: mass balance supplies the incompressibility
constraint, momentum balance evolves velocity under pressure, stress, and body
force terms, and boundary data specify how information enters through the domain
boundary
\cite[Part~I, Ch.~1, Secs.~2.2--3.1]{GaldiEtAl2008HemodynamicalFlows}
\cite[Sec.~6]{QuarteroniFormaggia2004CardiovascularModeling}.

\subsubsection{Conservation-law vocabulary}
\label{subsubsec:conservation-laws}

For a constant-density Newtonian lumen model, the familiar continuum tier can be
read schematically as
\begin{flalign*}
    \quad
    \Div\vu &= 0, &&\\
    \quad
    \rho\left(\pd_t\vu+(\vu\cdot\nabla)\vu\right)
    &= -\nabla p+\eta\Delta\vu+\rho\vfb. &&
\end{flalign*}
Generalized-Newtonian variants replace the constant viscosity by an effective
viscosity law. Those 3D equations are not solved in the main 1D realization, but
they explain why reduced models must declare which cross-sectional quantities are
retained and which transverse effects are closed.

The reduced model used later keeps an area-flow state. Cross-sectional averaging
turns local velocity and stress information into section area, flow, and source
or closure terms. As a result, the selected 1D equations are balance laws with
geometry, wall, profile, friction, rheology, and boundary choices built into the
model contract.

\subsubsection{Boundary-data vocabulary}
\label{subsubsec:boundary-initial-data-vocabulary}

For a single-vessel domain, the boundary is decomposed into inlet, outlet, and
wall portions. A 3D problem may prescribe velocity, flow, traction, pressure, or
wall-coupling data depending on the tier. A 1D problem inherits the same need for
a declared boundary rule, but the incoming information is applied to the reduced
state rather than to the full velocity field. This is why
Section~\ref{subsec:initial-boundary-conditions} treats the fixed-area
characteristic boundary rule as part of the implemented solver, and why
Section~\ref{subsec:boundary-diagnostics} records subcritical characteristic
signs for the comparison runs.

The continuum background stops at this vocabulary. The report does not use a
full 3D derivation as evidence for the 1D discretization; the relevant evidence
comes from the implemented balance law, manufactured-solution checks, the
geometry-rest test, and the declared velocity observation operator.
